\section{Skolem--Mahler--Lech theorem}

\begin{definition}[Arithmetic progression]
\label{def:sml-ap}
\lean{SkolemMahlerLechTheorem.arithProg}
\leanok
$\operatorname{arithProg}(a,d)=\{a+dn:n\in\mathbb N\}$.
\end{definition}

\begin{lemma}[$p$-adic zero-set dichotomy]
\label{lem:sml-padic}
\lean{SkolemMahlerLechTheorem.padic_zero_set_dichotomy}
\notready
The zero set of a linear recurrence eventually either vanishes identically or avoids zero on each residue class modulo a suitable $d$. This $p$-adic analytic step remains a \texttt{sorry}.
\end{lemma}

\begin{theorem}[Eventual periodicity of the zero set]
\label{thm:sml-periodic}
\lean{SkolemMahlerLechTheorem.zeroSet_eventuallyPeriodic}
\notready
\uses{lem:sml-padic}
The zero set of a linear recurrence over a characteristic-zero field is eventually periodic.
\end{theorem}

\begin{theorem}[Decomposition of an eventually periodic set]
\label{thm:sml-decomp}
\lean{SkolemMahlerLechTheorem.eventuallyPeriodic_decomp}
\leanok
\uses{def:sml-ap}
An eventually periodic subset of $\mathbb N$ is a finite exceptional set together with a finite union of arithmetic progressions.
\end{theorem}

\begin{theorem}[Skolem--Mahler--Lech]
\label{thm:skolem-mahler-lech}
\lean{SkolemMahlerLechTheorem.MainTheorem}
\notready
\uses{thm:sml-periodic,thm:sml-decomp}
The zero set of a linear recurrence over a characteristic-zero field is a finite union of arithmetic progressions together with a finite exceptional set.
\end{theorem}
